% Responsable: Persona 4. ~150 palabras. Cada aplicación con una referencia revisada por pares.
\section{Desarrollos recientes y otras aplicaciones}

Aplanar una superficie curva sin deformar las formas es un problema que
aparece en muchas áreas. En gráficos por computadora, los mapas conformes por
mínimos cuadrados se usan para pegar texturas sobre modelos 3D
\cite{levy2002}, y en neuroimagen la corteza cerebral se lleva de forma
conforme a una esfera o a un plano para comparar cerebros de distintas
personas \cite{gu2004,hurdal2004}. Para superficies dadas como mallas de triángulos
existe una teoría de mapeos conformes discretos: el empaquetamiento de
círculos converge al mapeo de Riemann \cite{rodin1987}, y la equivalencia
conforme de mallas permite aplanarlas con distorsión controlada
\cite{springborn2008}. En física e ingeniería, la transformación de Joukowski
convierte el flujo alrededor de un cilindro en el flujo alrededor de un perfil
alar, y la de Schwarz--Christoffel resuelve problemas de potencial
(electrostática, flujo) en regiones poligonales \cite{driscoll2002}. Los
métodos numéricos también siguen avanzando: los solucionadores
\emph{lightning}, basados en funciones racionales, resuelven la ecuación de
Laplace con muchas cifras correctas y en fracciones de segundo, aun en
regiones con esquinas \cite{gopal2019},
justo el tipo de región que complicó el ajuste polinomial en este trabajo.
